% New insertion audit. All prior sources, main text and state remain unchanged.
\subsection{Inserting the restricted convolution estimate into the actual finite atoms}
\label{sec:lichtman-hb-modulus-insertion-20261009}

\paragraph{Precise conclusion.}
Write $T$ for the even target, $x$ for the scale of the complete integer
being distributed, $U$ for one individual Heath--Brown variable, and
$\theta=7/32$. The new convolution margin is denoted by $\sigma$.
The conclusion here is the signed distribution estimate for the finite
atom classes specified below, using the restricted Proposition~5.2
derived in \texttt{lichtman-prop52-retained-phase-bridge-20261009.tex}.
It is not an estimate for an arbitrary coefficient supported on integers
that happen to admit factorizations. The classical BV, rough SW,
one/two-long-variable and standard HB inputs are identified explicitly.
No density integral, Chen weight, or numerical $K,c,b$ is recertified.

Fix a finite buffered policy, all its ancestor labels, a positive minimum
outer prime exponent $a_*>0$, and a fixed maximum number of original
prime variables. At each enhanced rough atom require either the full
and selected-square licence below, or the strict $W$ licence below.
The current positive product upper tree, ordinary lower-start trees,
and admitted typed table descendants satisfy these particular interfaces;
the numerical table alone is not their licence. Low ranks and classical
states have the separate complete-modulus bound of paragraph~6.
All finite meshes and buffers are fixed before $T$ tends to infinity.

\paragraph{1. The coefficient identity, before distribution.}
A labelled rough box has atom
\[
 c_J\,\mathbf1_{\rm squarefree}
  (\chi_{J_1}*\cdots*\chi_{J_r})*\beta_J,
 \qquad |\beta_J|\le1,
\]
where squarefreeness concerns the rough product, and beta has only small
prime factors. The scalar contains the original repeated-label factorial
once. For each Type II scale, partition the labelled slots into groups
of cardinalities $r_1,r_2,r_3$. Divide each group convolution by $r_i!$
and replace the scalar by $c_J\prod r_i!$. This is an exact convolution
identity. Unique prime factorization bounds each normalized group by one,
including repeated interval labels. Attach beta to group three; uniqueness
of the small/rough split preserves this bound. Unit groups are point
masses at one, not proposed SW variables.

First remove rough square collisions in the analytic remainder. If the
rough threshold is $T^\omega$ and total support is $D_*\le T^{\vartheta_*}$,
then, with fixed divisor orders after HB,
\[
 E_{\rm coll}\ll T^{o(1)}
 \sum_{\substack{d\le D_*\\p^2\mid d\text{ for some }p>T^\omega}}
 \left(T/d+1+T/\varphi(d)\right)
 \ll (T+D_*)T^{-\omega+o(1)}.
\]
Thus the three independent group sequences can be used without an
unlicensed cross-group squarefree restriction. The $+1$ also handles
the exact complete product $T$. Every modulus prime remains coprime to
$T$. Signs and complex atom scalars are left unchanged; the absolute
convolution remainder accepts them. There are finitely many scalars with
bounded total absolute mass, independently of $T$.

An assignment is constructed at the ordered upper endpoints and frozen
on the entire box. It may depend on the atom, $T$, and the HB length $U$,
but not on the individual outer integer or progression errors. This
distinguishes the construction from pointwise integer factorization.
For example, arbitrary coupled coefficient matrices on two disjoint
prime intervals can have unbounded rank; support membership alone
provides no bounded finite sum of separated group coefficients.

\paragraph{2. Actual selected gates and their factor inequalities.}
Use local logarithms, a licence exponent $1/2\le\vartheta<5/8$,
$D=x^{\vartheta-2\zeta}$, and
$x^\zeta\le A\le x^{1/3-\zeta/2}$, with $0<\zeta<10^{-5}$.
For an actual frozen total cap $T^\ell$, put $q=\log_xT$ and choose
$\vartheta=q\ell+2\zeta$, so $D=T^\ell$. Any fixed endpoint inflation
is included in $\ell$ first. The original branch slack, then sufficiently
small $\zeta$, ensures that this auxiliary $\vartheta$ is still below
its local licence. Thus no raw endpoint is silently equated with the
source's smaller support exponent $\vartheta-2\zeta$.
Let $\alpha\ge\log_x T$ and put
\[
 v=2\vartheta-1,\quad
 u=\min\left\{\frac{2-(2+2\theta)\vartheta}{3-4\theta},
                 \frac{2-2\vartheta-\alpha\theta}{3}\right\}.
\]
The complete endpoint vector must satisfy every full square. For a
selected subset $I$ of its marked slots, also require, separately for
every unselected slot $j$,
\[
 b_I+\sum_{i<j,\ i\notin I}t_i+2t_j\le\log_xD,
 \qquad b_I=\sum_{i\in I}t_i,\qquad v\le b_I\le u.       \tag{H1}
\]
For slots after the marked prefix this is exactly the full square.
The upper odd-cube history implies its own and the next even square;
the lower even-cube history, after its initial square, implies its own
and the next odd square. Marked initial squares must be checked in
their actual order, or by the existing merge lemma; no sign argument
supplies missing squares.

Writing $b_0=x^{b_I}$, if $C=D/(Ab_0)\ge1$, greedily place the remaining
whole slots in capacities $A,C$. At each step their square in (H1) is
at most the product of the two remaining capacities, so at least one
accepts the whole next slot. This gives $d=abc$, $a\le A$, $b=b_0$,
$c\le D/(Ab_0)$, and directly
\begin{align*}
 A^2bc^2&\le D^2/b_0\le x^{1-4\zeta},\\
 T^\theta Aab^5c^2&\le x^{\alpha\theta}D^2b_0^3
                         \le x^{2-4\zeta},\\
 A^{1+2\theta}ab^{5-2\theta}c^{2+2\theta}
 &\le D^{2+2\theta}b_0^{3-4\theta}
                         \le x^{2-(4+4\theta)\zeta}.       \tag{H2}
\end{align*}
If $C<1$, all slots are reassigned into $A,B=D/A$, not just the
unselected ones. Full square support gives the same greedy proof.
Here $1<B<b_0\le x^u$ and $c=1$. The first bound becomes
\[
 A^2b\le AD\le x^{23/24-5\zeta/2}\le x^{1-3\zeta},
\]
and the other two follow by replacing $b_0$ in (H2) by $B$.
This is the required below-one-capacity repair, not a fictitious
integer factor of size less than one.

For a prime box perform these calculations on its ordered upper-endpoint
vector, allowing real endpoint sizes in the same greedy inequalities.
The resulting \emph{group upper bounds} are then fixed for all tuples
in the box. In the selected case write $b_+$ for the endpoint product:
the first bound is $A^2 B_2B_3^2\le D^2/b_+$, using
$B_2\le b_+$ and $B_3\le D/(Ab_+)$. It is not necessary to
reassign slots using an individual tuple or to divide this bound by
that tuple's potentially smaller selected product. The compact gate
trim and final endpoint widths supply all strict endpoint assumptions.

The equivalent tent restriction is
\[
 \vartheta\le V_\alpha(b_I):=
 \min\{(1+b_I)/2,\ [2-(3-4\theta)b_I]/(2+2\theta),
                 1-(\alpha\theta+3b_I)/2\}.              \tag{H3}
\]
At $\alpha=1$ its middle line is dominated and gives the familiar
two-line tent. It must be retained when making the local-scale check.
Singleton-largest, leading-product, nonleading-pair and arbitrary
retained mask branches all use (H1), with every complement slot included.
For instance mask138 selects $\{2,4,8\}$ and has five complement tests
at slots $1,3,5,6,7$; mask129 selects $\{1,8\}$ and has six tests.
No arbitrary mask is licensed by its sum alone.

\paragraph{3. The strict $W$ alternative and ordinary root.}
Here one of $p_1,p_2$ has exponent at most $v$, the other has square
at most $x^{1-\vartheta-2\zeta}$, and $p_3\le x^{u-v}$.
Initialize the capacities $A,x^v,D/(Ax^v)$ with only $p_1,p_2$.
All capacities are at least one. If no later slot fails, their product
bounds give, respectively, exponents
\[
 1-4\zeta,\quad \alpha\theta+8\vartheta-3-4\zeta,\quad
 (8-6\theta)\vartheta+4\theta-3-(4+4\theta)\zeta.
\]
The last two are at most $2-4\zeta$ and $2-(4+4\theta)\zeta$
because $u\ge v$. At the first failed slot $p_k$, append it to the
second partial product. Its new value lies in $(x^v,x^u]$, while
$Abc_0<D$ follows from the failed first capacity and the full square
at $k$. Now the two remaining capacities are legitimate and (H2)
completes the proof. This explicitly uses the repaired initialization.

For the ordinary lower starts, $t_1\le h=19/125<\mu=25/128$.
The largest-singleton line and the strict $W$ rescue give
\[
 \Phi\ge\max\{(1+t_1)/2,\Lambda-3t_1/8\}\ge249/448,
 \qquad\Lambda=153/256.
\]
The two lines intersect at $25/224$. For the rescue the other initial
prime has square exponent at most $2\mu<1-\Lambda$; if its first
$W$ gate fails, the singleton already dominates. Thus a slightly
buffered ordinary root is licensed, including its untouched descendants.
For the other $u-v$ branch the slack at $W(t)=\Lambda-3t/8$ is explicitly
$33/17-(107/34)W(t)-t=525/8704+(49/272)t>0$; hence both third-slot
conditions have been checked, with compact trimming where a boundary
condition is approached.
This calculation does not grant a blanket level to larger marked primes;
those keep their own selected/$W$ label and physical ceiling.

\paragraph{4. From the medium HB slot to the four repaired inequalities.}
The groups in this paragraph are the original \emph{distributed prime}
groups; they are not the three groups of modulus slots in paragraph~1.
Fix the original chart, of product scale $X=\prod_i R_i$, before
Mellin separation and HB. Write $X_+$ for its upper complete-product
endpoint. With $f(t)=t^{1/(\log\log t)^3}$ choose once
$Z=f(C_kX_+)$, where the fixed $C_k$ bounds the enlargement after all
at most $6k$ HB slots and regrouping/dyadic division of the complement.
Thus every later source scale $x=UV$ is within fixed multiplicative
constants of $X$ and is at most $C_kX_+$. The roughness condition is
automatic on the original primes. It is imposed before HB and never
changed on a later rectangle. The uniform rough-SW window
$x^{c_1/(\log\log x)^3}\le Z\le x^{c_2/(\log\log x)^3}$,
with fixed positive $c_1,c_2$, supplies this fixed $Z$.
In particular choose $C_k\ge2^{6k_*+2}$; then one may take
$(c_1,c_2)=(1,2)$ eventually. The full proof, including arbitrary
auxiliary integers, is
\texttt{lichtman-hb-rough-sw-insertion-20261009.tex}, paragraphs~2--6.

Apply order-three HB separately at each original prime scale.
A retained group has at most six slots and lower-endpoint product
at least its original prime scale/64. Select the medium class using
the \emph{original} chart:
\[
 X^{3\zeta}\le U\le X^{1/3+\zeta/4}.
\]
The two power gaps absorb all the preceding fixed constants. On every
eventual source rectangle this guarantees, for sufficiently large $T$,
\[
 x^{2\zeta}\le U\le x^{1/3+\zeta/2},\qquad A=Ux^{-\zeta}.
\]
Its coefficient is rough $1$ or rough $\mu$, times an interval,
normalized logarithm and bounded-variation smooth twist. The explicit
rough-SW proof supplies the complete auxiliary coprimality hypothesis;
an arbitrary grouped coefficient is not substituted for this slot.
Use the endpoint construction above with $\delta=\zeta$.
The weaker common savings before replacing $A$ by $U$ are
$3\zeta,3\zeta,2\zeta$. The replacement costs
$2\zeta,\zeta,(1+2\theta)\zeta$, leaving
\[
 \zeta,\quad2\zeta,\quad9\zeta/16.                       \tag{H4}
\]
Reserve $\zeta/16$ for local-coordinate effects and another
$\zeta/16$ for beta and fine endpoints. Take $0<\sigma\le\zeta/1000$.
The least remaining saving $7\zeta/16$ exceeds $15\sigma$, so
\begin{align*}
 Q_1&<Ux^{-\sigma},& U^2Q_2Q_3^2&<x^{1-8\sigma},\\
 T^\theta UQ_1Q_2^5Q_3^2&<x^{2-15\sigma},&
 U^{1+2\theta}Q_1Q_2^{5-2\theta}Q_3^{2+2\theta}
   &<x^{2-15\sigma}.                                   \tag{H5}
\end{align*}
Beta enlarges only group three: its costs are $2\ell,2\ell,
(2+2\theta)\ell$, which fit the separate reserve. Smaller dyadic
factor lengths preserve all four bounds since $5-2\theta=73/16>0$.
The complement of the medium slot is divisor-bounded. Truncate by
a divisor-function logarithmic threshold and normalize, using the new
divisor-coefficient repair. The discarded progression tail uses the
nonzero moment of $\tau_C(n)\tau_C(|T-n|)$ and treats $n=T$ separately
by $T^{\vartheta_*+o(1)}$; the coprime mean has only logarithmic loss.
The selected slot already is one-bounded after logarithm normalization.
The roughness threshold, original group constraints and SW constants
are uniform over all these rectangles. A bound valid only for a
newly chosen threshold on each rectangle would not suffice here.

\paragraph{5. Moving target, with the new smaller residue range.}
For ordinary marked counts, $x$ is the scale of the original prime,
not $T/b$ or the residual sieve level. For switched counts it is the
scale of the complete product, not a selected quotient. The retained
scales satisfy $T^{1-\kappa_0}/C_0\le x\le C_1T$, with fixed constants.
Consequently $T\ge x/C_1$, supplying the previously missing lower
residue bound in the repaired proposition.

Choose $\kappa=\sigma/200$ and, for example,
$0<\kappa_0<\sigma/(1000+\sigma)$. Then
$\kappa_0/(1-\kappa_0)<\sigma/1000<\kappa$ and eventually
$T<x^{1+\kappa}$, with $\kappa\le\sigma/100$.
These choices replace the weaker old condition involving one common
epsilon. Shrink them further if required by a fixed support margin.
In particular, with $q=\log_xT$, selected tent gaps under rescaling
are $q$ times their old gaps minus, respectively,
\[
 (q-1)/2,\quad(q-1)/(1+\theta),\quad(q-1)
\]
when $\alpha=q$. Using a common larger $\alpha$ costs an additional
$\theta(\alpha-q)/2$ in the last line. Full and complement-square
gaps scale homogeneously. For the two $W$ third-slot constraints the
corresponding subtractive constants are $33(q-1)/17$ and
$5(q-1)/3$, plus the small $\alpha-q$ term. The designated-prime
gate $t_i\le2\vartheta-1$ improves by $q-1$ when $q\ge1$.
The \emph{other} initial-prime square gap instead becomes
$q(1-\vartheta-2t_j)-(q-1)$ before the support buffer; it needs its
own positive reserve. If $q<1$ because $x>T$, then $1-q=O(1/\log T)$
and the same fixed reserves absorb that case.
Thus the same finite positive margin survives, rather than silently
using the $\alpha=1$ endpoint on smaller scales.

All these choices are after the common compact trim and before $T$.
The beta length, rough cutoff, bounded rank and final fine width are
then chosen in that order as in the existing compiler. The signed tail
below $2T^{1-\kappa_0}$ costs $T^{1-\kappa_0+o(1)}$ after divisor
summation; it imposes no positive lower bound on $\kappa_0$.
No error of size $x^{o(1)}$ is multiplied into a merely logarithmic
remainder; fixed logarithmic losses are paid by requesting a stronger
logarithmic exponent from the repaired proposition.

\paragraph{6. The low-rank alternative really is classical.}
Before beta, endpoint and local-coordinate losses, the available gaps are:
\[
 \begin{array}{c|c|c}
 \text{family}&\text{rough support exponent}&1/2-\text{exponent}\\\hline
 \text{upper product tree}&2(19/96)=19/48&5/48\\
 \text{ordinary lower starts}&2h=38/125&49/250\\
 \text{enhanced marked/table families}&(\Lambda+57/224)/2=1527/3584
                                      &265/3584
 \end{array}
\]
Choose the total losses smaller than, for example, half the minimum
of these raw gaps. The resulting complete supports still have fixed
positive gaps below $1/2$; the displayed exact gaps are not claimed
unchanged after those losses.
The last row uses the actual full square $t_1+2t_2\le\Lambda$ and
$t_1<57/224$, so $t_1+t_2\le(\Lambda+t_1)/2$. It does not use the
smaller upper-tree first cutoff on those marked states. Descendants
append smaller primes. Classical states already fix total level below
$1/2$. The sufficiently small local-scale change preserves each gap.

For one original prime, classical BV bounds the absolute modulus sum
with a maximum over all primitive residue classes. For several original
prime factors, select one prime interval as the SW factor and combine
the rest into a divisor-bounded factor, both at fixed positive power
scales. Motohashi's generalized BV below $1/2$ has that same maximum;
the precise statement is Polymath v4, Claim2.6 and Theorem2.8. Complex
coefficients follow by linear decomposition. Thus $T\bmod d$ is included
for every $(T,d)=1$, even when $T$ varies. Bounded compiled coefficients
multiply this absolute estimate harmlessly. No enhanced L23 theorem,
fixed-residue assumption, or conjectural GEH beyond $1/2$ enters here.

\paragraph{7. No-medium alternatives and the exact scope of insertion.}
If no medium slot exists in the buffered original chart, every slot
is below $X^{3\zeta}$ or above $X^{1/3+\zeta/4}$. Every original group
has a long slot; otherwise its product is at most
$X^{18\zeta+o(1)}$, contradicting its retained original scale when
$18\zeta<a_*/2$. Three long slots have product exceeding
$X^{1+3\zeta/4}$, impossible since the retained full dyadic product
is $O_k(X)$. There are at most two long slots and at most two original
groups. Neither long slot is a truncated Mobius slot.
For a uniform conservative parameter write $\zeta_L=2\zeta$.
The short product is at most $x^{10\zeta_L+o(1)}$ in the one-long
case (at most five shorts, giving $15\zeta+o(1)$), and
$x^{20\zeta_L+o(1)}$ in the two-long case (at most ten shorts,
giving $30\zeta+o(1)$). One long slot therefore has length at least
$x^{1-10\zeta_L-o(1)}$ and is treated by the elementary rough progression
count (L23 Lemma3.4), including its exact coprime-mean correction.
With two, use the same fixed $Z$ and choose
\[
\vartheta_* /(1-\kappa_0)<3/5
 <\tfrac23(1-20\zeta_L)-4\sigma,\qquad
 1-10\zeta_L>3/5+2\sigma.
\]
The rough divisor-function estimate follows from Maynard II v1
Lemma7.1 and its Proposition5.2 proof: the former is pointwise in every
primitive residue, and the latter explicitly allows individual interval
endpoints. Here the stronger sufficient modulus bound is intentional:
the printed proposition and its last proof line do not have identical
dependence on the short factor. The explicit reconstruction in
\texttt{lichtman-hb-long-branches-20261009.txt}, paragraphs~3--6,
retains the frequency-cutoff losses, extends the roughness proof to
this chart-fixed $Z$, and controls interval boundaries by the nonzero
shifted moment with the complete product $T$ isolated. Thus it does
not import the old fixed-shift smoothing argument unchanged.
The underlying two-Poisson/Weil bound has a power saving, so small divisor
convolutions from the fundamental lemma cost only $x^{o(1)}$ there.
Partial summation allows the required smooth twists, not arbitrary
arithmetic weights on the long variables.

These alternatives and (H5) give, for the same signed finite coefficient,
\[
 \sum_d\rho_T(d)\left\{
 \sum_{p_1\cdots p_k\equiv T\ (d)}W_T
 -\frac1{\varphi(d)}\sum_{(p_1\cdots p_k,d)=1}W_T\right\}
 =O_A(T\log^{-A}T).
\]
The constants depend on the fixed policy, cutoff and finite smooth
family, not the even target. This is an absolute signed-error assertion;
a positive product upper bound still needs its separate positive
majorant and actual density proof. No arbitrary signed atom has been
reclassified as a positive counting object.

The explicit scope consists of licensed selected/full-square and strict
$W$ atoms, their finite typed descendants, and the classical low-rank
alternatives above. The extra two-large-marked-prime candidate has a
separate three-case partition with six strict monomial margins (minimum
$117/31250$); it also fits after taking smaller parameters, but is not
silently inserted into the current ledger. No claim is made for a
high-rank atom with only carrier level $\Theta<3/5$, an arbitrary outer
residue filter, a vanishing cutoff, or a mesh growing with $T$.
Within the specified classes no new incompatible hypothesis of the
restricted proposition remains. End-to-end numerical/density/pointwise
assembly and the reported Wu issue are outside this insertion audit.

\paragraph{Source ledger and the meaning of the checks.}
The complete atom identity and collision calculation were compared with
\texttt{actual-k-distribution-joint-audit.tex}; the parity/full-square
translation with \texttt{own-lower-compiler-chain-audit-20261009.tex};
the marked/table low-rank bound with
\texttt{own-chen-table-insertion-chain-20261009.tex}; and the actual
eighth masks and carrier/cutoff distinction with
\texttt{current-eighth-distribution-interface-audit-20261009.tex}.
These files supply the displayed structural hypotheses, not a substitute
for the repaired large-residue estimate. The selected and $W$
calculations above rederive the requisite arithmetic rather than
quote their old analytic status. Primary text anchors are L23 v1
Definition~2.1, Proposition~5.2 and Lemmas~6.3--6.4; Maynard II v1
Lemma~7.1 and Proposition~5.2; and Polymath v4 Claim~2.6/Theorem~2.8.
Their precise extracted-line ranges are recorded in the companion JSON.
The rational checker verifies exponent identities, margins and parameter
compatibility; it is not a numerical replay of the analytic theorems.
